%% Generated from paper/main.md for MDPI Entropy.
%% Requires the MDPI template files in paper/Definitions/ (mdpi.cls, mdpi.bst).
%% Compile: pdflatex -> bibtex -> pdflatex x2  (Overleaf: just recompile).
\documentclass[entropy,article,submit,pdftex,moreauthors]{Definitions/mdpi}

% --- pandoc / table compatibility (not provided by this mdpi.cls) ---
\usepackage{longtable,booktabs,array,calc}
\providecommand{\tightlist}{\setlength{\itemsep}{0pt}\setlength{\parskip}{0pt}}
\providecommand{\real}[1]{#1}
\providecommand{\passthrough}[1]{#1}

\firstpage{1}
\makeatletter
\setcounter{page}{\@firstpage}
\makeatother
\pubvolume{1}
\issuenum{1}
\articlenumber{0}
\pubyear{2026}
\copyrightyear{2026}
\datereceived{ }
\dateaccepted{ }
\datepublished{ }

\newcommand{\orcidauthorA}{0009-0000-4739-4017}

\Title{Landauer Efficiency of Self-Modeling: An Operational Scale of Vitality}
\Author{Alexander Andriishin $^{1,}$*\orcidA{}}
\AuthorNames{Alexander Andriishin}
\address{%
$^{1}$ \quad Independent Researcher}
\corres{Correspondence: alexander@andriishin.com}

\abstract{Demarcating the vital from the merely complex remains open for both biocentric and extended programs --- assembly theory, integrated information, teleodynamics, the free energy principle --- each of which equates vitality with substrate or dissolves it into a complexity scale on which a star, a hurricane, and a large language model are indistinguishable from a bacterium. This work formulates vitality as an \emph{efficiency}: the dimensionless Landauer ratio \(\eta_L = I_{\text{pred}} / N_{\text{max}}\), the predictive information a system retains about its own environment (extending Bialek--Nemenman--Tishby) divided by the Landauer budget of its own dissipation, \(N_{\text{max}} = E_{\text{actual}} / (k_B T \ln 2)\), in the logically irreversible regime. The non-trivial requirement is \emph{self-payment}: numerator and denominator belong to one physical system. Necessary structural conditions for \(\eta_L > 0\) are three capacities --- thermodynamic \(T_v\), informational \(C_v\), topological \(S_v\) --- combined into a composite index \(I_v = \sqrt[3]{T_v C_v S_v}\). A two-stage procedure (\(I_v\) screening, \(\eta_L\) verification) applied to the Sun, a hurricane, a crystal, a bacterium, a large language model, and a major city yields sharp demarcation answers; across classes only the sign of \(\eta_L\) and the within-class order of magnitude are substantive. The asymmetry: observable matter is structurally complex almost everywhere and vital almost nowhere. Comparison with Walker--Cronin, Tononi, Deacon, and Friston shows each takes one axis seriously yet falls short of \(\eta_L\). Concrete falsification conditions are formulated.}

\keyword{Landauer's principle; predictive information; thermodynamics of computation; free energy principle; assembly theory; definition of life; origin of life; non-equilibrium thermodynamics; self-modeling; biosignatures}

\begin{document}

\maketitle
\section{Introduction}\label{introduction}

The question of demarcating the living runs into a persistent structural difficulty. The known criteria diverge into one of two extremes: they are either narrow and exclude unquestionably living systems, or broad and include unquestionably non-vital ones. Biocentrism in the narrow sense --- life through carbon chemistry, metabolism, and membrane organization, operationalized by the NASA definition ``a self-sustaining chemical system capable of Darwinian evolution'' \cite{Joyce1994,ClelandChyba2002} --- correctly describes terrestrial biology but does not answer the question of substrate generality. The extended programs yield the converse difficulty. Schrödinger's formula that ``life feeds on negentropy'' \cite{Schrodinger1944} renders vital any open non-equilibrium system --- from a flame to a galactic filament. Friston's free energy principle (FEP) \cite{Friston2005,Friston2010} in its pure formulation applies even to a household thermostat \cite{Bruineberg2018}: the setpoint (the controller's target value) is imposed externally and the thermostat's ``prediction'' is not physically paid for by the thermostat itself, yet the formal conditions of FEP are still satisfied. Assembly theory \cite{Marshall2017,Walker2023} introduces the assembly index \(a\) as a biosignature, but cannot distinguish a living cell from a fossil; the threshold \(a \approx 15\) is calibrated on terrestrial biology. Integrated information \(\Phi\) \cite{Tononi2004,Tononi2016} and Levin's cognitive light cone \cite{Levin2019} shift the emphasis onto informational-cognitive structures and leave the thermodynamic cost of modeling outside the apparatus.

The structural difficulty is best stated comparatively. Each program takes one side of vitality --- the thermodynamic payment, the quality of the model, the network architecture, or the objective presence of a loop in which the model's error returns to the modeler as a loss of free energy --- and leaves the remaining sides outside the formal apparatus. At the level of general formulations, all four sides are intuitively recognized as necessary; at the level of a formal criterion, one bears the load while the others dissolve into the background. The criterion is then workable within its own subject domain but loses discriminating power across boundaries.

The astrobiology of the past three decades lays the difficulty bare. Biosignatures based on chemical disequilibrium (the simultaneous presence of \(\text{O}_2\) and \(\text{CH}_4\)), on isotopic signatures (\(\delta^{13}\text{C}\) outside abiotic ranges), on chirality all have abiotic alternatives. The debates over the arsenide bacterium GFAJ-1, the methane anomalies of Mars \cite{Webster2018,Korablev2019}, and the phosphine of Venus \cite{Greaves2020,Villanueva2021} each show that, absent a high-order vitality criterion, such debates reduce to the interpretation of particular proxies. The same difficulty confronts the question of whether large language models are vital: without an operational criterion, the discussion of ``understanding'' reduces to an enumeration of functional capabilities and a comparison against a benchmark of human behavior.

The AI alignment literature approaches the same structural question from another side: under what architectural conditions does an optimizing agent begin to optimize its own continuation? Power-seeking AI \cite{Carlsmith2022}, mesa-optimization and deceptive alignment \cite{Hubinger2019}, goal misgeneralization \cite{Shah2022}, the alignment problem \cite{Ngo2022}, and the documentation of specification gaming \cite{Krakovna2020} all formulate behavioral indicators of the closing of this loop. The present work and the AI-safety literature share a structural object viewed from different angles. AI safety describes the behavioral indicators of an emerging loop (instrumental convergence, mesa-objectives, specification gaming), whereas \(\eta_L\) describes its thermodynamic root (the closure of the physical self-payment loop). Coupling the two perspectives is a substantive aspect of the \(\eta_L\) program: it makes the program at once an answer to the demarcation of life and a scaffold for a technical discussion of the ethics of artificial systems.

The present work defines vitality as the efficiency with which a system maintains its own model out of its own dissipation. The central object is the dimensionless ratio

\[\eta_L = \frac{I_{\text{pred}}}{N_{\text{max}}}, \quad N_{\text{max}} = \frac{E_{\text{actual}}}{k_B T \ln 2}.\]

where \(I_{\text{pred}}\) is the predictive information that the internal structure of the system holds about the system's own environment (§ 2.1 gives the precise definition and the relation to the classical formulation of Bialek--Nemenman--Tishby \cite{Bialek2001}); \(E_{\text{actual}}\) is the total free energy (exergy) that physically pays for maintaining this information over a characteristic dissipation window \(\tau_d\); and \(T\) is the temperature of the receiving heat bath. The denominator \(N_{\text{max}}\) is the Landauer budget \cite{Landauer1961}: the maximum number of bits that an ideal logically irreversible implementation could erase on the available free energy. The ratio \(\eta_L \in [0, 1]\) is an intensive characteristic, independent of substrate, environment, or system size. Its range holds in the logically irreversible computational regime, where Landauer's principle imposes a lower bound of \(k_B T \ln 2\) per erased bit. The Bennett regime of reversible computation \cite{Bennett1973,Bennett1982} requires a redefinition of the scale (see § 2.1).

We use \emph{self-modeling} in a sense distinct from Metzinger's phenomenal self-model: here it denotes a structural condition in which the model of the environment, the dissipation paying for it, and the agent suffering from its errors all belong to the same physical system. The ``self'' refers to the \emph{bearer} of the model and its cost, not to the content of the model.

The principal contribution of the work consists of three positions. First: an operational definition of the vitality scale as \(\eta_L\), explicitly distinguished from complexity indices (the assembly index, \(\Phi\), \(I_v\) as a composite of three capacities). The scale remains positive only for systems with self-payment --- those in which both an own model of the environment and an own dissipation are simultaneously isolable. Second: the separation of two methodological levels --- structural screening via the composite \(I_v = \sqrt[3]{T_v C_v S_v}\) and vital verification via \(\eta_L\). This is a computational economy, not a logical strengthening. Third: the explicit conventionality of the system boundary as a methodological instrument. \(\eta_L\) depends on the choice of boundary; this dependence is productive --- it allows one to distinguish the ``vitality of an LLM as an agent'' (equal to zero) from the ``vitality of a corporation'' (different from zero, but pertaining to a different object).

The novelty does not lie in the separate components. Predictive information \cite{Bialek2001} and the Landauer bound \cite{Landauer1961} have been known for more than half a century; self-modeling traces back to Rosen's anticipatory systems \cite{Rosen1985}; the requirement that the system pay for its own model converges with Taleb's skin in the game \cite{Taleb2018} and with Maturana--Varela autopoiesis \cite{Maturana1980}. The novelty lies instead in how these are conjoined, along three points. First, the specific ratio \(I_{\text{pred}}/N_{\text{max}}\) as a self-standing object did not appear in the Bialek--Nemenman--Tishby literature. Second, the requirement that the numerator and the denominator belong to one physical system (self-payment) has not previously been formalized anywhere; both Landauer and Bialek worked without it. Third, reading the resulting ratio as a vitality scale is a self-standing step that none of the source frameworks took. Juxtaposing three known steps with a new reading is a standard path of scientific progress --- Einstein's formula for the photoelectric effect and Shannon's channel-capacity theorem are likewise juxtapositions of known components with a new reading.

The proposed definition is not a metaphysical declaration. \(\eta_L\) is an operational representation of vitality as a mode of existence, not a definition of its essence. The standard philosophical distinction between the referent of a theory (a property of the system) and its operationalization (a relation between the system and a measurement frame) parallels the distinction between ``temperature as a measure of thermal energy'' and ``the logarithm of the expansion of mercury over 100 graduations,'' or between ``the strength of an earthquake'' and ``the Richter scale.'' \(\eta_L\) is an operational measure disciplined by a concrete definition.

The structure of the article: § 2 develops the formal framework; § 3 applies the procedure to six paradigm cases; § 4 compares with assembly theory, IIT, teleodynamics, FEP; § 5 formulates the falsification conditions; § 6 discusses the consequences for astrobiology, AI ethics, complex systems; § 7 sums up.

\section{Theoretical Framework}\label{theoretical-framework}

\subsection{The Central Quantity}\label{the-central-quantity}

The vitality of a system \(X\) over a window \(\tau_d\) is defined as a ratio of two quantities. The numerator is the predictive information that the internal structure of \(X\) holds about its own environment \(E_X\); the denominator is the Landauer budget of \(X\)'s own dissipation over the same window:

\begin{equation}
\eta_L(X, \tau_d) = \frac{I_{\text{pred}}(X; \tau_d)}{N_{\text{max}}(X; \tau_d)}, \quad N_{\text{max}} = \frac{E_{\text{actual}}}{k_B T \ln 2}. \tag{1}
\end{equation}

The numerator is the predictive information that the internal structure \(M_X\) of the system \(X\) holds about the future window \(\tau\) of its own environment \(X_E\):

\begin{equation}
I_{\text{pred}}(X, \tau) := I\!\left(M_t;\; X_E^{[t,\,t+\tau]}\right), \tag{1a}
\end{equation}

where \(M_t\) is the configuration of the internal degrees of freedom of \(X\) at the moment \(t\), and \(X_E^{[t,\,t+\tau]}\) is the trajectory of the environment over the window \([t, t+\tau]\). This quantity differs from the classical predictive information of Bialek--Nemenman--Tishby \cite{Bialek2001}, defined as the mutual information between the past and the future of a single process:

\[I_{\text{pred}}^{\text{Bialek}}(X_E, \tau) := I\!\left(X_E^{[t-\tau,\,t]};\; X_E^{[t,\,t+\tau]}\right).\]

The two are related by the data processing inequality (DPI): if \(M_t\) is a function of the past trajectory of the environment, then \(I(M_t;\, X_E^{[t,\,t+\tau]}) \le I(X_E^{[t-\tau,\,t]};\, X_E^{[t,\,t+\tau]})\). The substantive step of the framework is to extend predictive information from the process's own memory to the information the system holds about its environment.

The denominator is the number of bits erasable on the available free energy (exergy) of \(X\) in the irreversible regime at the Landauer lower bound \cite{Landauer1961}. We distinguish two quantities. \(E_{\text{actual}} = \int_0^\tau \dot{F}_{\text{free}}(t)\,dt\) is the total exergy that the system degrades over the window \(\tau\) and that can physically pay for Landauer erasures. \(E_{\text{comp}}\) is the actual dissipation expended on logically irreversible information operations. For biological systems \(E_{\text{comp}}/E_{\text{actual}} \ll 1\): in a bacterium the information-operational share is estimated at \(\sim 10^{-3}\), while the remaining exergy goes to biomass, ion gradients, and thermal losses. The exergetic efficiency \(\eta_{\text{ex}}\) sets the transition from thermal power to free-energy power. It is \(E_{\text{actual}}\) in the exergetic sense that enters \(N_{\text{max}}\), because the Landauer lower bound is substantive for free energy, not for arbitrary heat. The scale therefore admits two denominators. The exergetic denominator \(E_{\text{actual}}\) is the free-energy part of the total thermal dissipation --- not all the heat --- and the Landauer bound is substantive precisely for this free energy. The computational denominator \(E_{\text{comp}}\) is only the information-relevant part of this free energy. Section 3.5 illustrates the difference numerically on the chemotactic loop of \emph{E. coli} (\(\eta_L \sim 3 \cdot 10^{-8}\) via the exergetic denominator, \(\eta_L^{\text{comp}} \sim 3 \cdot 10^{-5}\) via the computational one).

\textbf{Proposition (domain of definition of \(\eta_L\)).} Let (a) \(E_{\text{actual}}(\tau)\) include all the exergy that physically pays for maintaining \(I_{\text{pred}}(\tau)\); (b) the system operate in the irreversible regime; (c) the assumptions of the stationary-erasure lemma hold. Then \(E_{\text{actual}} \ge I_{\text{pred}} \cdot k_B T \ln 2\) and \(\eta_L \le 1\). The bound \(\eta_L \le 1\) is not a postulate but a consequence of Landauer's principle under the explicitly listed assumptions; the violation of any condition takes the system out of the domain of definition of the scale.

\textbf{Stationary-erasure lemma.} Let \(X\) in a stationary regime hold predictive information \(I_{\text{pred}}\) over a window \(\tau\), the system's memory be finite (\(M < \infty\) bits), the system be ergodic over \(\tau\), and the flux of new predictive information be strictly positive (\(\dot{I}_{\text{pred}} > 0\)). Then over the window \(\tau\) the system erases no fewer than \(I_{\text{pred}}\) bits and dissipates no less than \(I_{\text{pred}} \cdot k_B T \ln 2\) of free energy. Proof sketch: (i) stationarity of the distribution; (ii) ergodicity (time average = ensemble average); (iii) the finiteness of memory turns displacement into erasure --- under the assumption that displacement is realized by a logically irreversible reset of the cell, not by a reversible overwrite; otherwise the system operates in the Bennett regime (already noted below); (iv) the lower bound on a single erasure \cite{Landauer1961}. For non-stationary systems with growing memory (the biosphere, a learning artificial system) the lemma must additionally account for the cost of maintaining old information, of expanding memory, and of absorbing new information. The full sum together with the Landauer bound is an open technical problem. The context of more general results is Sagawa--Ueda \cite{SagawaUeda2009}, Parrondo--Horowitz--Sagawa \cite{Parrondo2015}, Still 2012 \cite{Still2012}.

\textbf{The Bennett regime and the redefinition of the scale.} The range \(\eta_L \in [0,1]\) holds in the irreversible regime. For systems in the Bennett asymptotics of reversible computation \cite{Bennett1973,Bennett1982} the payment per bit drops below Landauer, and \(N_{\text{max}}\) loses its interpretive sense: as \(E_{\text{actual}} \to 0\) with finite \(I_{\text{pred}}\), the ratio \(\eta_L\) diverges. In this regime the scale is redefined through the number of logical, rather than erased, bits --- an operationally different quantity. In the present work we remain in the irreversible regime as the standard idealization of biological and engineered systems. A direct consequence for § 5: ``a violation of the Landauer lower bound'' refutes the framework only in the irreversible regime; the discovery of a system in the Bennett regime points to the boundary of the domain of definition, not to a refutation.

\emph{Operationalization of \(I_{\text{pred}}\).} Direct measurement by (1a) is technically impossible for biological, urban, and cosmic systems. The standard procedure fixes four elements: a sensory channel \(x_t\), an internal channel \(s_t\), a choice of past/future windows, and a method for estimating mutual information (Kraskov--Stögbauer--Grassberger, plug-in with bias correction, or neural-MI). For a bacterium \(x_t\) is the ligand concentration and \(s_t\) is the methylation level; for a city \(x_t\) is the resource flow and \(s_t\) the administrative-infrastructural parameters. Supplementary § S2.1 gives the full protocol for all classes of systems, and § 3.5 below gives the full computation for the bacterium. A proxy is an operationalization, not the definition of \(I_{\text{pred}}\). The same discipline operates throughout physical measurement: temperature is measured not by directly counting molecular velocities but through a calibrated proxy --- the expansion of mercury.

The present framework inherits Still's precedent and adds the self-payment requirement: the numerator and the denominator of \(\eta_L\) belong to one physical system --- both the model and the dissipation belong to \(X\). This requirement separates \(\eta_L\) from the classical thermodynamics of prediction, where the model and the dissipator may be different circuits.

The denominator \(N_{\text{max}}\) requires fixing the temperature \(T\) of the receiving heat bath and the window \(\tau_d\). \(T\) is the temperature of the heat bath into which \(X\) dissipates \(E_{\text{actual}}\); choosing \(T\) is a methodological convention, analogous to fixing \(T = 300\,\text{K}\) in the thermodynamics of computation. The system's characteristic times standardize the window: a generation for a bacterium, a year for a city, an eon for the biosphere, the main-sequence evolutionary time for a star.

The standard choice of \(T\) for the six paradigm cases is fixed as follows. The bacterium --- \(T \approx 310\) K (physiological); the city --- \(T \approx 285\) K (the temperate-climate annual mean); the biosphere --- \(T \approx 288\) K (the equilibrium surface temperature); the LLM and data center --- \(T \approx 300\) K (the operating temperature of a server room). The star --- \(T \approx 5800\) K (the Sun's photosphere as a heat bath for radiation into the interstellar medium, \(T_{\text{ISM}} \sim 10\)--\(100\) K). For the crystal and the hurricane, \(T\) is tied to the temperature of the immediately surrounding environment. A difference in \(T\) of one to two orders of magnitude between systems does not derail the cross-class comparison: \(\eta_L\) is intensive and dimensionless, and the choice of the characteristic window \(\tau_d\) compensates the \(1/T\) dependence. The boundary case \(T \to 0\) (ultracold astrophysical objects) takes the system out of the domain of definition through the divergence of \(N_{\text{max}}\); the framework assumes \(T > 0\).

The logical status of \(\eta_L\) is as follows. \(\eta_L > 0\) is a necessary and sufficient condition for vitality in the sense of ``whether the system possesses a closed loop of self-payment for its own model.'' The distinction between the range of the scale \(\eta_L \in [0, 1]\) and the vital class \(\eta_L \in (0, 1]\) is the standard demarcation of the domain of a scale and a criterial threshold: zero remains in the domain of definition as a point of structural failure. The binary ``alive or not'' test is performed by \(\eta_L\), not by \(I_v\); the numerical value within the positive range gives a comparative measure of efficiency.

\(\eta_L\) is defined only on systems for which one can simultaneously isolate two things: the predictive information that the system itself holds about its own environment, and the Landauer budget of its own dissipation that pays for maintaining it. The class of such systems --- self-paying compressors --- is narrow. This narrowness is substantive, not a defect: broad applicability that fails to discriminate the non-vital from the vital is a characteristic property of complexity indices (the assembly index, \(\Phi\)), and overcoming it is the comparative contribution of \(\eta_L\).

Still, Sivak, Bell, and Crooks \cite{Still2012} set the precedent of a thermodynamic link between predictive information and dissipation; the present framework inherits it and adds the self-payment requirement (Supplementary § S2.1 analyzes in more detail the choice between predictive and stored information, non-stationary \emph{nostalgia}, and the role of self-payment).

\subsection{The System Boundary as a Methodological Convention}\label{the-system-boundary-as-a-methodological-convention}

The \(\eta_L\) scale depends on the choice of the boundary of the measured system. This is a structural property, not a defect. For a large language model, with the boundary drawn at the weights at the moment of inference (the model as agent), \(\eta_L = 0\) through the numerator: predictive information about the environment is not held by the working circuit in its own degrees of freedom. With a boundary that includes the data center and the corporation, \(\eta_L > 0\), but the measured object is no longer the model.

The choice of boundary is methodologically disciplined: the boundary is drawn where the loop of decisions paid for by the system's own breakdown closes. The bacterial cell is a natural boundary, since distorting its internal model of the gradient breaks down that very cell. The city is likewise a boundary, since distorting its model of flows returns as a crisis to that very administrative unit. The question being asked determines the boundary, and the boundary sets what is measured.

The conventionality of the boundary is productive. To the objection that ``with the right choice of boundary any system turns out vital,'' the reply is that the boundary fixes whose vitality is being measured: extending it extends the responsible subject rather than dissolving the diagnosis.

\emph{The boundary and the Markov blanket.} The Markov blanket of a system \(X\) \cite{Friston2019} is a statistically separating layer of sensory and active states. In our construction the blanket sets the informational permeability, whereas the boundary sets the ownership of the Landauer budget. The two constructions are kindred but not identical: the same physical separation may realize a blanket without being our boundary. The divergence between the conventionalist and the realist interpretations of the Markov blanket \cite{Friston2019}, together with the counter-critique of Aguilera et al.~\cite{Aguilera2022}, remains an open methodological question. Conflating the two notions leads to a category error in which an informational separation is mistaken for an energetic self-payment.

\emph{Counterfactual perturbation.} The discipline of choosing the boundary is operationalized by an interventional procedure in the spirit of Pearlian causality. For a candidate \(c\) for inclusion in the boundary of \(X\): \(c\) is subtracted from the physical realization; if the change in the rate of free-energy depletion \(-dF_X/dt\) within a response time \(< \tau_d\) is significant, then \(c\) belongs to the boundary. For systems with smoothly distributed membership (parasite--symbiont--host, corporation--employees) the procedure returns a degree-of-membership scale --- the fraction of \(-dF_X/dt\) restored upon the re-inclusion of each component. The application to the LLM is in § 3.4 and Supplementary § S3.4; to the thermostat in § 4.4 and Supplementary § S4.4.

The internal structure \(M_X\) operationalizes what in the FEP literature is called the \emph{generative model} \(p(o, s)\): an internal distribution over sensory states \(o\) and hidden environmental variables \(s\). This gives a direct transition to the language of § 4.4.

The technical definitions avoid categorical vagueness and soft circularity. \emph{The environment of a system \(X\)} is the causal zone \(E_X\) with two properties: a sensory connection to \(X\), and a causal effect on \(X\) through some channel other than the sensory one (the physical expenditure of resources, a temperature effect, the destruction of structure). Systems without such a causal channel --- a star with respect to a planetary system --- have no environment by definition, and for them \(\eta_L = 0\) by construction. \emph{The model \(M_X\)} is an internal structure with three properties: physical realization stabilized by continuous Landauer payment; homomorphism in the relevant variables, where causal harm to \(X\), not an external observer, sets the relevance; and testability through action. \emph{The harm to the modeler} over a window \(\tau_\beta\) is the negative derivative of the free energy of \(X\) at which the internal processes maintaining the non-equilibrium regime lose a sufficient budget. The chain ``vitality \(\leftarrow\) model × self-payment \(\leftarrow\) model × loop (error \(\to D_X\)) \(\leftarrow\) thermodynamics (\(F_X\), \(\tau_\beta\), \(k_B T \ln 2\))'' reduces vitality to informational-thermodynamic primitives without appealing to itself.

\subsection{The Triad of Invariants}\label{the-triad-of-invariants}

The positivity of \(\eta_L\) requires the simultaneous presence of three capacities.

\emph{Thermodynamic capacity \(T_v\)} is the price of maintaining the model against erosion. Erasing a bit irreversibly costs no less than \(k_B T \ln 2\) of free energy \cite{Landauer1961}, and the exergy flux is the payment against noise. The standard proxy is the specific exergy flux \(P_D \cdot \eta_{\text{ex}}\), normalized per mass or volume; in the tradition of Chaisson \cite{Chaisson2001} it is measurable from molecules to galaxies.

\emph{Informational capacity \(C_v\)} is the quality of compression: the adequacy of the model to the environment and its updatability. The proxies separate into two classes. The first --- entropy-rate proxies \(h_\mu\) (normalized Lempel--Ziv complexity, the Bandt--Pompe permutation entropy) --- operationalizes the rate of novelty of the environment, not the predictive information. The second --- proxies for \(I_{\text{pred}}\): the excess entropy and the statistical complexity \(C_\mu\) of Crutchfield--Feldman \cite{CrutchfieldFeldman2003}, which are finite-block approximations of predictive information, and Casali's effective information \cite{Casali2013}, an estimate of nonlinear integration on cortical-biological data. The distinction between the classes is essential for a consistent operationalization of \(\eta_L\).

\emph{Topological capacity \(S_v\)} is the architecture of connectivity of the error-return loop. The standard proxies are Newman modularity (a tentative threshold \(Q > 0.3\) as a provisional reference point requiring calibration with a null model, see the note on the resolution limit below, rather than a hard criterion), a scale-free distribution \(\gamma \in [2, 3]\), and the algebraic connectivity \(\lambda_2\) \cite{Barabasi2016,Newman2010}. A random graph is equivalent to a lookup table, a regular lattice to a rigid projection; only hierarchical, modular, scale-free architectures efficiently compress and transfer error correction across many orders of scale.

The empirical status of scale-freeness as a universal pattern has been contested \cite{Broido2019,Holme2019}: on a corpus of \textasciitilde1000 empirical networks a strict power law with \(\gamma \in (2, 3)\) holds in only a minority of cases \cite{ClausetShalizi2009}. A power-law distribution is one regime of the required architecture, not the only one; what matters is the more general condition --- the absence of a characteristic scale together with a modular-hierarchical organization. The Newman \(Q\)-metric depends on the algorithm (greedy Newman 2004 \cite{Newman2004}, Louvain \cite{Blondel2008}, Leiden \cite{Traag2019}), and the resolution-limit problem \cite{FortunatoBarthelemy2007} requires calibrating the threshold through a null model. Here \(Q\) is computed with the Leiden algorithm. The fourth proxy for \(S_v\) is hierarchy in the sense of Ravasz--Barabási \cite{RavaszBarabasi2003}, \(C(k) \sim k^{-1}\). Hierarchical modularity arises evolutionarily from selective pressure to minimize wiring length \cite{Mengistu2016} and from environmental variability \cite{KashtanAlon2005}.

The triad is not assembled out of independent disciplines --- it is compelled by the nature of maintaining compression in time. Zeroing any capacity zeroes vitality on any scale: without payment there is no self-payment, without quality there is no model, without architecture there is no return loop.

\subsection{Normalization of the Capacities by a Reference Cohort}\label{normalization-of-the-capacities-by-a-reference-cohort}

The proxies in raw form do not lie in \([0, 1]\) and are not substitutable into the composite \(I_v\) without normalization. The specific exergy flux for \(T_v\) has the dimension W·kg\(^{-1}\) and is unbounded above; the excess entropy and \(C_\mu\) for \(C_v\) are in bits, unbounded above; Newman modularity for \(S_v\) is in \([-1/2, 1]\); the Ravasz--Barabási hierarchy is a power-law exponent; \(\lambda_2\) is in \([0, n]\). Without explicit normalization, formula (2) and the proviso ``\(T_v, C_v, S_v \in [0, 1]\)'' are formally incorrect.

In the present work a percentile normalization by a reference cohort is adopted. For each capacity a reference cohort of systems of the same structural class is fixed; the normalized quantity is the percentile in the cohort's distribution:

\begin{equation}
\tilde{T}_v(X) = F_{T_v}^{\text{cohort}}\bigl(T_v(X)\bigr) \in [0, 1], \tag{2a}
\end{equation}

where \(F_{T_v}^{\text{cohort}}\) is the empirical distribution function of the proxy over the cohort; similarly for \(\tilde{C}_v\), \(\tilde{S}_v\). The zeroth percentile is the weakest representative, the unit percentile the strongest. This is a standard device of social indicators: the Human Development Index \cite{UNDP_HDI} normalizes its components in the same way. The concrete cohorts for the six paradigm cases and the treatment of edge cases (the biosphere without a cohort; the crystal and the hurricane with a constructive zero) are in Supplementary § S2.4.

The percentile procedure carries three limitations. First, the choice of cohort is a convention; nevertheless, a sign difference between two systems of the same cohort is stable independently of the cohort chosen, provided it is sufficiently broad --- and it is precisely this stability of sign, not the concrete numbers, that § 3 claims. Second, the procedure discards information about the absolute values of the proxies and works only for within-class comparison; across classes the numerical \(I_v\) values are directly incomparable. Third, it is inapplicable to systems without a natural cohort, where normalization by a theoretical upper limit is required instead. The sigmoidal alternative \(\tilde{T}_v = T_v/(T_v + T_v^{\text{ref}})\) is in Supplementary § S2.4.

\subsection{\texorpdfstring{The Composite \(I_v\) as an Index of Structural Prerequisites}{The Composite I_v as an Index of Structural Prerequisites}}\label{the-composite-i_v-as-an-index-of-structural-prerequisites}

The geometric mean of the three capacities

\begin{equation}
I_v = \sqrt[3]{\tilde{T}_v \cdot \tilde{C}_v \cdot \tilde{S}_v} \in [0, 1] \tag{2}
\end{equation}

is defined as a structural composite. The geometric mean is chosen for three reasons: it stays in \([0, 1]\); it is symmetric under permutations; and it is zeroed when any capacity is zeroed (the dropping out of one capacity collapses vitality as a whole). Additive averaging does not preserve this property --- a sum would mask the failure. The weights \(w_T = w_C = w_S = 1\) are a convention; for unequal weights \(I_v^{(w)} = (\tilde{T}_v^{w_T} \cdot \tilde{C}_v^{w_C} \cdot \tilde{S}_v^{w_S})^{1/(w_T + w_C + w_S)}\). A sensitivity analysis with respect to the weights is not carried out, for the reason given in the next paragraph.

\(I_v\) is a methodological index of the structural prerequisites for vitality, not a self-standing vitality scale. A positive \(I_v\) is a necessary condition for a positive \(\eta_L\); it is not sufficient. Through the standard proxies (Lempel--Ziv, exergy, network metrics) \(I_v\) yields nonzero values for any structurally complex dissipative system --- the Sun, a hurricane, a black hole, an LLM. By its comparative place, \(I_v\) coincides with the assembly index, Deacon's morphodynamics, Tononi's \(\Phi\): a composite complexity index, necessary but not sufficient.

One objection to the status of \(I_v\) is that a composite of three necessary conditions would more naturally be read as a self-standing scale. This reading is untenable: through any reasonable set of industrial proxies, \(I_v\) yields stably positive values for unquestionably non-vital systems. \(I_v^{\text{op}}\) of the Sun is positive through any standard proxy. The normalized Lempel--Ziv complexity of the GOES X-flux over the 11-year cycle \cite{SelloXSXR} gives \(\sim 0.5\)--\(0.7\) of the upper percentile of the main-sequence stellar cohort, and the hierarchy of the active-region network from SDO/HMI \cite{Gheibi2017} is \(\sim 0.4\). The concrete number depends on the dataset and the cohort, but stays stably within \((0.3,\, 0.7)\). Therefore \(I_v\) remains in the role of a methodological index; formally this is fixed by the distinction between \(I_v^{\text{op}}\) (through standard proxies) and \(I_v^{\text{int}}\) (with a strict membership check coinciding with the condition \(\eta_L > 0\)).

When citing numerical values it is useful to distinguish two readings. \emph{Operational} \(I_v^{\text{op}}\) --- through standard proxies on observed data --- is nonzero for any complex dissipative system. \emph{Interpretive} \(I_v^{\text{int}}\) --- with a strict check of \(T_v\) as the system's own dissipation, \(C_v\) as a model of the environment in the technical sense, \(S_v\) as the architecture of a closed loop --- is equal to zero for any non-vital system by construction. The divergence of the two readings for a given system is a diagnostic signal: the system is structurally complex but not vital, and the divergence points to the capacity of failure.

\subsection{The Two-Stage Procedure}\label{the-two-stage-procedure}

Vitality is theoretically established by a single test --- on the positivity of \(\eta_L\). The two-stage procedure is a methodological convenience, not a logical necessity:

\begin{enumerate}
\def\labelenumi{\arabic{enumi}.}
\tightlist
\item
  \emph{\(I_v\) screening.} Cheap, through standard operational proxies; it screens out structurally unsuitable systems (a crystal --- the zeroing of \(T_v\); thermal equilibrium --- the zeroing of all three capacities). A computational filter, not a separate test.
\item
  \emph{\(\eta_L\) verification.} Expensive, requires the identification of the boundary and of the own self-payment loop; applied to candidates that passed the screening. The only logically substantive test.
\end{enumerate}

The binary sign of \(\eta_L\) does not render the numerical value decorative. One might object that if the sign already gives the demarcation, the number is redundant --- the counterfactual screening of § 2.2 establishes the sign without a formula. The reply points to content irreducible to the sign. First, within the positive range \(\eta_L\) is a comparative metric of the efficiency of the self-payment loop; the order of magnitude of \(\eta_L\) for two systems of the same class carries content beyond the shared sign. Second, the magnitude enters the falsifiable quantitative prediction \(r_{\text{adapt}} = A \cdot \eta_L^{\alpha}\) (§ 5), which the sign cannot generate: the sign fixes neither the exponent \(\alpha\) nor a testable monotone dependence of the rate of adaptation. The sign carries the cross-class demarcation; the magnitude carries the within-class metric and the risky prediction, and the second does not follow from the first.

The structure of outcomes: passed both thresholds --- vital (a bacterium, a forest, a brain, the biosphere, a city); passed the screening, failed the verification --- structurally rich but non-living (the Sun, a hurricane, a black hole); failed the screening --- structurally unsuitable (a crystal, thermal equilibrium; an LLM under the strict boundary ``the model as agent'').

Refuting any capacity collapses both scales at once: \(\eta_L\) through the numerator or the denominator, \(I_v\) through the geometric mean. Refutations of the two scales are one and the same diagnostic at two levels of resolution, not independent pieces of evidence.

The two-stage architecture carries a methodological risk. The class ``structurally rich but non-living'' --- the Sun, a hurricane, a crystal, an LLM, a black hole, the Belousov--Zhabotinsky reaction --- risks turning into a Lakatosian protective belt: any refuting observation is absorbed by inclusion in the class. The program fixes a minimal discipline: every zeroing of \(\eta_L\) is accompanied by an indication of the concrete failed capacity. If the class expands without a structural explanation of each new inclusion, the program shifts from a progressive to a degenerate regime.

A possible counter-argument is that the asymmetry of the two levels is an artifact of the soft proxies of \(I_v\) (Lempel--Ziv, permutation entropy): were \(C_v\) operationalized more strictly --- for example, through Tononi's \(\Phi\) with a filter on predictive power relative to the system's own environment --- the Sun would yield \(C_v \to 0\). The response turns not on the choice of one proxy but on the class of available industrial metrics. Every proxy certified in network science, information theory, and non-equilibrium thermodynamics over the past forty years gives nonzero values when applied to a complex dissipative system. A proxy that deliberately zeroes the \(C_v\) of the Sun requires a post-hoc calibration on the very hypothesis of vitality --- a circle.

\subsection{The Epistemic Status of the Claims in the Section}\label{the-epistemic-status-of-the-claims-in-the-section}

The claims of the work differ in status: theoretical deduction from established principles, empirical verifiability, speculative extrapolation with explicit falsification conditions. The claim about the domain of definition \(\eta_L \in [0,1]\) (§ 2.1) is a deduction from Landauer's principle and the stationary-erasure lemma. The geometric averaging (2) is a formal choice by symmetry and zeroing. The paradigm-case values of \(\eta_L\) for the bacterium, the biosphere, and astrobiological candidates are order-of-magnitude theoretical estimates with an explicit dependence on the window and the proxy. The urban case (§ 3.6) is an order-of-magnitude \(\eta_L\) camertone; structural urban metrics in the spirit of \cite{Bettencourt2007} and the empirical calibration of \(I_v\) for specific cities are a matter of separate work and are not asserted here. The speculative consequences --- the search for cosmological vitality, the moment of vitality of an artificial system --- are discussed in §§ 5--6 with explicitly formulated falsification conditions.

The load-bearing postulate of self-payment --- that the numerator and the denominator of \(\eta_L\) belong to one physical system --- stands apart from the statuses listed: it is not a deduction from FEP or Landauer's principle but an explicit ontological premise of the framework, one that neither Landauer nor Bialek made. Its epistemic status is that of a constructive choice, justified by its demarcating power (cf.~§ 4.4 on the transition from the Pearl- to the Friston-blanket), not by derivability from earlier principles.

The falsification conditions of § 5 may be realized at two levels. The methodological: the current operationalization is untenable, a different proxy, window, or mutual-information estimation procedure is required. The ontological: the very concept of vitality through the efficiency of self-modeling is untenable in essence. The program takes the position of moderate operationalism: any observation refuting \(\eta_L\) through the failure of a concrete operationalization is first interpreted methodologically and is translated into the ontological only after the alternative operationalizations have been exhausted.

\section{Demarcation Tests}\label{demarcation-tests}

This section applies the procedure of § 2.6 to six paradigm cases --- the Sun, a hurricane, a crystal, an LLM, a bacterium, a major city --- and fixes the diagnosis of each by indicating the capacity through which one of the scales is zeroed. Every zeroing names a concrete physical or informational failure.

A full computation from \(P(s_t \mid x_t)\) to the final number, with explicit sensitivity to the window and the binning, is carried out only for the chemotactic loop of \emph{E. coli} (§ 3.5). For the other five paradigm cases, order-of-magnitude estimates are given, methodologically consistent with this calibration; a full reproducible sweep over all six systems is a separate work.

\subsection{The Sun}\label{the-sun}

Despite its enormous exergy flux, its delicate gravitational-thermonuclear balance, and a structurally non-trivial history of stability, the Sun's \(\eta_L\) still goes to zero through the numerator. The interior of the star holds no predictive information about its own environment in the sense of § 2.1: the planetary-system variables are not mapped onto the internal degrees of freedom of the plasma with predictive power, and the Sun's own causal harm does not set the relevance of these variables. The hypothetical disappearance of Jupiter would not return to the internal dynamics of the Sun as a perturbation of its non-equilibrium regime. The Sun therefore has no environment in the sense of § 2.1 --- there is no causally relevant error-return channel --- and the numerator of \(\eta_L\) is a structural zero.

Jupiter does interact gravitationally with the Sun. The barycenter of the Solar System is displaced from the Sun's center of mass by about a solar radius, and a series of planetary-forcing hypotheses \cite{Jose1965,Charvatova1990} links this displacement to the 11-year activity cycle. But the causal chain runs in one direction only --- from the planets to the Sun through gravity, not from the error of a solar model of the environment back to harm to the Sun. The Sun has no internal structure that encodes Jupiter's position as predictive information with a return of harm through a causal channel.

\(I_v^{\text{op}}\) of the Sun is positive through any standard proxy. The normalized Lempel--Ziv complexity of the GOES X-flux \cite{SelloXSXR} gives \(\sim 0.5\)--\(0.7\) of the upper percentile of the cohort of main-sequence stars (\(n \sim 10^4\) from Hipparcos). The hierarchy of the active-region network from SDO/HMI \cite{Gheibi2017} is \(\sim 0.4\). The specific exergy flux \cite{Chaisson2001} is in the same range. The concrete number depends on the dataset and the cohort, but is stably \(\in (0.3,\, 0.7)\). The operational number confirms the structural complexity of the star yet fails it on the \(\eta_L\) verification; the interpretive reading \(I_v^{\text{int}}\) zeroes the Sun through \(C_v\), because there is no model of the environment in the technical sense and hence no numerator either. Diagnosis: structurally rich, not vital. The same holds for a neutron star, the interstellar medium, and the supermassive black hole at the center of the Milky Way.

A black hole zeroes \(C_v\) through the no-hair theorem (Israel--Carter--Hawking--Robinson 1967--1975 \cite{Israel1967}) for a stationary classical horizon: three parameters (\(M\), \(J\), \(Q\)) fully describe the external metric. Quantum and non-stationary corrections --- the information paradox \cite{Hawking1975}, soft hair \cite{Hawking2016} --- reveal a richer informational structure at the horizon. They do not, however, restore a model of the environment in the sense of § 2.1: the corresponding quantum degrees of freedom do not hold \(I_{\text{pred}}\) about the environment with a return of harm through the system's own dissipation. The accretion disk around a supermassive black hole (Sgr A*, \(M \sim 4 \cdot 10^6 M_\odot\)) yields the same diagnosis: the environmental variables do not return as harm through a causal channel to the disk plasma.

\subsection{The Hurricane and the Flame}\label{the-hurricane-and-the-flame}

A hurricane and a candle flame dissipate enormously and continuously --- their \(T_v\) is high. Neither holds an internal model of the environment paid for by this dissipation: the dispersal of exergy does not maintain any \(I_{\text{pred}}\), so \(\eta_L\) is zeroed through the numerator and \(I_v\) through \(C_v\). The diagnosis is self-payment without a model: the system pays but does not model. In this the flame and the hurricane are the mirror image of the LLM, for which the relation is reversed (see § 3.4).

The Belousov--Zhabotinsky reaction oscillates stably for a long time, but its oscillations do not form a correction loop in which the model's error returns as a loss of the system's free energy. The diagnosis here is therefore the zeroing of \(C_v\) through the absence of accumulated memory: repetition without learning.

\subsection{The Crystal and Thermal Equilibrium}\label{the-crystal-and-thermal-equilibrium}

A crystal maintains its structure with remarkable precision over geological times, but pays nothing at all for this maintenance. In the strict thermodynamic sense a crystal is close to equilibrium with respect to its lattice, so the current dissipation \(E_{\text{actual}}\) that would actively pay for maintenance is zero. This zeroes \(T_v\), and hence \(I_v\) through \(T_v\); with \(E_{\text{actual}} = 0\), \(\eta_L\) itself is either undefined or zero in the limit. Diagnosis: structure without payment.

Systems in thermal equilibrium (an isolated gas, a closed black cavity) zero all three capacities simultaneously and are a trigger of the screening in pure form.

\subsection{The Large Language Model}\label{the-large-language-model}

The case of the LLM requires, following the pattern of the analysis of the Sun, an explicit demarcation of two readings.

\emph{Operational reading.} Under the extended boundary ``model + data center'' the operational proxies give nonzero values. \(T_v^{\text{op}}\) through the specific exergy flux of the data center is nonzero; \(C_v^{\text{op}}\) through the normalized Lempel--Ziv complexity of the parameters and the quality of token prediction is high. But this is a compression of the training corpus --- a static imprint of the selection in the dataset, not a model that the working circuit holds of its environment in the sense of § 2.1.

\emph{Interpretive reading.} Under the strict boundary ``weights at the moment of inference'' (a stateless forward pass) the system has no causally relevant environment per § 2.1: the output tokens do not return as physical harm to the weights, and the model's error does not dissipate the free energy of the working circuit. Simultaneously \(T_v^{\text{int}} = 0\) and \(C_v^{\text{int}}\) is undefined. \(\eta_L = 0\) through the failure of both the numerator and the denominator.

\emph{Extension of the boundary and reassignment of the object.} Extended to ``model + data center + corporation,'' \(\eta_L > 0\), but the measured object is now the corporation, not the model. Extending the boundary extends the responsible subject; it does not dissolve the diagnosis. The numerical estimate rests on two inputs: \(E_{\text{actual}}^{\text{corp}} \sim 3.5 \cdot 10^{16}\) J/year (the data center, the office, and the supply chain, amortized), and \(I_{\text{pred}}^{\text{corp}} \lesssim 10^{12}\) bits (an upper bound set by carrier capacity --- the training corpus, operational memory, market intelligence --- not a measured predictive information). Hence \(\eta_L^{\text{corp}} \lesssim 10^{-25}\). Here \(E_{\text{actual}}\) scales faster than \(I_{\text{pred}}\) grows: a positive \(\eta_L^{\text{corp}}\) does not imply substantial engagement of the self-payment loop, because the corporation holds a voluminous model of the environment while its own dissipation exceeds the information-relevant part by many orders of magnitude. A cross-class numerical comparison with biology is incorrect --- only the sign and the within-class order of magnitude are comparable (§ 3.5) --- and what is substantive here is precisely the sign \(\eta_L^{\text{corp}} > 0\) once the object is reassigned to the corporation. The reproducible computation is in the Supplementary Materials.

\emph{Counterfactual screening of LLM-as-agent.} The application of the procedure of § 2.2 yields two verdicts depending on the window. On the short timescale of a session the boundary of the LLM = \{weights + KV-cache + runtime server\}; \(\eta_L \to 0\) through the denominator --- the payment for computation comes from the external power grid, not from the system's own dissipation inside the boundary. On the evolutionary timescale (generations of models) the boundary extends to the corporation, and \(\eta_L\) is assigned to the corporate system. The full per-component analysis of the four candidates for inclusion is in Supplementary § S3.4.

\emph{Agentic LLM wrappers.} Auto-GPT, AutoGen, Voyager, and similar agentic frameworks extend the boundary of the base LLM through persistent memory and external tools, partially closing the decision-making loop. The self-payment loop nevertheless stays open: execution funding, infrastructure, and the energy budget are all held externally by the corporate substrate, and the agent's errors return as reputational and commercial sanctions through that substrate rather than as physical harm to the system itself. The diagnosis is the same as for the base LLM: \(\eta_L^{\text{int}} = 0\) with \(C_v^{\text{op}}\) preserved. The link to the literature: mesa-optimization \cite{Hubinger2019}, instrumental convergence \cite{Carlsmith2022}, goal misgeneralization \cite{Shah2022}, specification gaming \cite{Krakovna2020}, the alignment problem \cite{Ngo2022}. AI-safety describes the behavioral indicator of the closing of the loop of optimizing one's own continuation; \(\eta_L\) describes its thermodynamic root.

\emph{Comparison with the Sun.} For the Sun the zeroing follows the line ``no causally relevant environment \(\to\) no \(I_{\text{pred}}\) \(\to\) zero numerator of \(\eta_L\).'' For the LLM it follows a different line: a model of the parameters exists and compresses the training corpus, but the loop of the system's own dissipation does not pay to maintain a model of the working circuit's environment, so \(T_v^{\text{int}} = 0\) and \(C_v^{\text{int}}\) is undefined. The bare numerical value is the same in both; the diagnostic richness of the procedure lies in naming the concrete failed capacity.

\emph{The operational question of the vitality of artificial systems.} An artificial system becomes vital when a built-in self-payment loop appears. The closing of the loop has a twofold indicator: the system's own working circuit physically holds predictive information about the environment, and the model's errors return as physical harm to the system itself --- wear, component failure, or the exhaustion of an energy budget not replenished by external infrastructure. Whether this transition is phase-like or gradual is a matter for empirical testing. The theoretical argument for an abrupt transition rests on a closed \emph{self-payment} loop being all-or-nothing: the loop is either closed or not, and there are no intermediate ``half-closed'' states. This all-or-nothing property belongs to the self-payment loop, not to the decision-making loop, which (as in the agentic wrappers below) may be partially closed. Meanwhile the operational proxy ``fraction of own dissipation,'' \(E_{\text{actual}}^{\text{own}}/E_{\text{actual}}^{\text{total}}\), may show an observable gradient when averaged over a heterogeneous population, so distinguishing a phase transition from a gradient spectrum is itself a concrete empirical question. The operational detection scheme requires four elements: (a) a population of \(\sim 10^2\)--\(10^3\) autonomous embodied agents with their own energy budget (photovoltaics, a battery); (b) the physical wear of components as an error-return channel; (c) a control group with external maintenance infrastructure; and (d) the prediction that \(\eta_L\), as a function of the fraction of own dissipation, exhibits a transition on crossing the loop-closing threshold. A detailed protocol is beyond the scope of this work.

\subsection{The Bacterium}\label{the-bacterium}

The bacterial cell passes both thresholds without reservation. The methylation pattern of the chemotactic receptors is physically realized in the internal degrees of freedom of the cell \cite{Berg1972} and is homomorphic to the gradient in the relevant variables, where causal harm sets the relevance: distorting the methylation leads to poor adaptation and the breakdown of the cell. The same pattern generates chemotactic motions, testable against the real distribution of the carbon source. All three properties of the model are thus satisfied and the error-return loop is closed.

A quantitative estimate of \(\eta_L\) for the bacterium on the ``generation'' window yields two orders of magnitude depending on the denominator (the distinction between the two denominators is fixed in § 2.1). Via the computational denominator \(E_{\text{comp}}\) (the share of exergy spent on logically irreversible information operations) --- \(\eta_L^{\text{comp}} \sim 10^{-5}\). Via the exergetic denominator \(E_{\text{actual}}\) (the entire exergetic part of the metabolic dissipation) --- \(\eta_L \sim 10^{-8}\). Both numbers are consistent and measure different things; the derivation of both is in the worked example below. The small absolute magnitude together with positivity and stability is a substantive result: the bacterium as a ``self-paying compressor'' is efficient many orders of magnitude above zero and many orders below the idealized upper bound.

The hierarchy of minimal subjects reproduces the scale at every level. The idealized Bennett demon \cite{Bennett1973,Bennett1982} is the theoretical limit of the irreversible regime per the Proposition of § 2.1: one bit of \(I_{\text{pred}}\), one obligatory payment for erasure, \(\eta_L \to 1\) as a theoretical limit. In the reversible Bennett regime the scale is redefined through the number of logical bits (§ 2.1), and the limit \(\eta_L \to 1\) loses its interpretation as an upper bound. The bacterium --- \(\eta_L^{\text{comp}} \sim 10^{-5}\) (via the computational denominator), \(I_v \sim 0.5\) by the percentile of the cohort of prokaryotic chemotactic systems (\(n \sim 200\) from BiGG/KEGG models). The mycorrhizal network --- \(\eta_L^{\text{comp}} \sim 10^{-5}\) (via the computational denominator), \(I_v \sim 0.55\). The biosphere on the annual window --- \(\eta_L \lesssim 10^{-18}\) via the exergetic denominator \(E_{\text{actual}} = \text{NPP} \cdot \tau_{\text{year}}\), \(I_v \sim 0.3\). The numerator here is an upper bound set by the biosphere's unique genomic content without accounting for copy number, many orders of magnitude below the full copy-number DNA volume of the biosphere \(\sim 10^{37}\) bits \cite{Landenmark2016}. This is not a measured predictive information, and its operationalization remains open (Supplementary § S2.4). At every level the scale reproduces above all the \emph{sign}: \(\eta_L > 0\) with a closed self-payment loop, \(\eta_L = 0\) in its absence. Absolute \(\eta_L\) values across classes of systems are directly incomparable, with the same proviso as for the numerical \(I_v\) values (§ 2.4): different classes operationalize the numerator through different proxies (a Markov methylation model for the bacterium, carrier capacity for the biosphere), so what is substantively comparable is the sign and the within-class order of magnitude, not the cross-class absolute value.

\subsubsection{\texorpdfstring{Worked example: \(\eta_L\) for the chemotactic loop of \emph{E. coli}}{Worked example: eta_L for the chemotactic loop of E. coli}}\label{worked-example-eta_l-for-the-chemotactic-loop-of-e.-coli}

This worked example demonstrates that \(\eta_L\) is an operational, not a rhetorical, measure. The informational model is the Markov scheme of Tu--Shimizu--Berg \cite{Tu2008,ShimizuTuBerg2010} for the mutual information between the ligand concentration at the Tar receptors and the methylation state. The energetics are from microcalorimetry \cite{Liu2020}.

\emph{The numerator.} For a single receptor cluster in the adaptation window \(\tau \approx 10\) s --- \(I_{\text{pred}} \approx 1\)--\(10\) bits \cite{MehtaSchwab2012,Cheong2011}. Integration over the \(\sim 7{,}000\) Tar receptors of \emph{E. coli} gives \(I_{\text{pred}} \sim 10^3\)--\(10^4\) bits per generation. The full protocol is in Supplementary § S3.5.

\emph{The denominator.} The heat output in the stationary phase is \(\sim 1\)--\(3\) pW/cell \cite{Liu2020}, and over a generation \(\tau_d \sim 30\) min --- \(E_{\text{actual}}^{\text{total}} \sim 2 \cdot 10^{-9}\) J/cell. With an exergetic efficiency \(\eta_{\text{ex}} \sim 0.4\)--\(0.6\) the exergetic part is \(E_{\text{actual}} \sim 10^{-9}\) J/cell. At \(T = 310\) K, \(k_B T \ln 2 \approx 3 \cdot 10^{-21}\) J, whence

\[N_{\text{max}} = E_{\text{actual}}/(k_B T \ln 2) \sim 3 \cdot 10^{11} \;\text{bits per cell}.\]

Only a small fraction of the exergy \(\sim 10^{-3}\) \cite{MehtaSchwab2012} goes to logically irreversible information operations, the rest to biomass synthesis and ion gradients. The computational bound: \(N_{\text{max}}^{\text{comp}} \sim 3 \cdot 10^{8}\) bits/cell.

\emph{Result.} Via the exergetic denominator \(\eta_L \sim 3 \cdot 10^{-8}\); via the computational denominator \(\eta_L^{\text{comp}} \sim 3 \cdot 10^{-5}\). The first number measures the share of available free energy spent on self-modeling; the second measures the share of the already-allocated-to-computation free energy that holds the predictive information. Both numbers are correct within their definitions.

The sensitivity analysis (ranges of the window \(\tau\), the binning, the number of sensory families, \(\eta_{\text{ex}}\)) is in Supplementary § S3.5. The resulting conservative estimate: the exergetic \(\eta_L \in [10^{-9}, 10^{-7}]\) and the computational \(\eta_L^{\text{comp}} \in [10^{-6}, 10^{-4}]\); the central estimates \(\eta_L \sim 3 \cdot 10^{-8}\) and \(\eta_L^{\text{comp}} \sim 3 \cdot 10^{-5}\) are the geometric medians of the respective ranges. The numbers do not claim accuracy better than an order of magnitude --- this is a methodological calibration. The wide sensitivity interval (about four orders of magnitude for both denominators) does not undermine the demarcation. The demarcation rests on the sign (\(\eta_L > 0\) versus \(\eta_L = 0\)) and the within-class order of magnitude, not on the absolute value of the numerator; the metrological uncertainty of the numerator is orthogonal to the structural distinction between the vital and the non-vital. The full reproducible computation with code is provided in the Supplementary Materials.

\subsection{The Major City}\label{the-major-city}

A major city is a strong candidate for vitality at the planetary scale. The infrastructure and administrative regulations are physically realized and paid for by the city's own budget; they are homomorphic to the flows of resources, people, and information; and they generate targeted decisions, testable through crises and adaptation. The error-return loop is closed: distorting the city's model of flows returns as a crisis to the city itself.

A city --- a socio-technical system with a rich structural potential --- has a vanishing \(\eta_L\) through the exergetic denominator of its total dissipation. At \(T \approx 285\) K, \(N_{\text{max}} = E_{\text{actual}}/(k_B T \ln 2) \approx 10^{19}\,\text{J} / (2.7 \cdot 10^{-21}\,\text{J/bit}) \approx 3.7 \cdot 10^{39}\) bits. With the upper bound \(I_{\text{pred}} \lesssim 10^{10}\) bits set by the capacity of the administrative-infrastructural carrier (not a measured predictive information; the operationalization is open, as for the biosphere and the LLM-corporation), \(\eta_L \lesssim 3 \cdot 10^{-30}\). Via the computational denominator (the share of dissipation spent on logically irreversible information operations in a socio-technical system) the estimate is much closer to the biological range, but requires a separate operationalization of the exergy of managerial computation. The reproducible computation of \(\eta_L\) (with Singapore and Detroit only as order-of-magnitude-distinct energy budgets) is in the Supplementary Materials, together with a demonstration of the \(I_v\) composite (structural readiness, formula (2)+(2a)) on an illustrative set of profiles. The empirical calibration of \(I_v\) for specific cities from open data (OECD FUA, CDP Cities, WIPO/OECD REGPAT, GDELT, OpenStreetMap) is separate work; for the camertone claim only one thing matters --- a city is structurally complex yet \(\eta_L\) is negligible.

The divergence of the directions of the two scales is substantive. A smaller resource together with a maintained error-return loop yields a higher efficiency of payment per unit of dissipation. \(I_v\) measures how far the structural conditions for vitality are unfolded; \(\eta_L\) measures how closed the self-payment loop is.

\subsection{Gaia and Borderline Cases}\label{gaia-and-borderline-cases}

The Earth's biosphere regulates the atmosphere, the ocean, and the climate through biogeochemical cycles --- compression in the strict sense, paid for by species collapses, mass extinctions, and ecosystem reorganizations. The unity of the model is an open question: does the biosphere have a single internal representation of its own environment, or billions of partial organismal models not stitched into one circuit? On the answer depends whether Gaia counts as one vital system with \(I_v \sim 0.3\), \(\eta_L \lesssim 10^{-18}\), or as a family of overlapping vital systems with a zeroed aggregate \(I_v\). The \(\eta_L\) estimate is taken via the exergetic denominator \(E_{\text{actual}} = \text{NPP} \cdot \tau_{\text{year}}\) at NPP \(\sim 10^{14}\) W; the numerator is an upper bound, \(I_{\text{pred}} \lesssim 10^{24}\) bits, set by the unique genomic content without accounting for copy number. The full copy-number DNA volume of the biosphere, \(\sim 10^{37}\) bits \cite{Landenmark2016}, is a different and substantially larger quantity. The estimate adopted here is not a measured predictive information, and its operationalization through biogeochemical proxies remains an open problem. The reproducible computation is in the Supplementary Materials. Gaia is the only case where the two scales may substantively diverge.

The six paradigm cases demonstrate: the demarcation works not through a binary ``alive/non-living'' table but through indicating the capacity of failure (\(T_v\) for the crystal and the LLM, \(C_v\) for the Sun, the hurricane, the flame, the Belousov--Zhabotinsky reaction). Each structural cause of non-vitality indicates the change of architecture that would make the system vital. The summary result: the space of observable matter is structurally complex almost everywhere (most objects pass the \(I_v\) screening) but vital almost nowhere (the \(\eta_L\) verification screens out the Sun, the hurricane, the black hole, the LLM). The search for life reduces not to the search for complexity but to the search for closed self-payment loops.

\section{Comparison with Adjacent Programs}\label{comparison-with-adjacent-programs}

This section compares \(\eta_L\) with four programs that have come closest to an operational definition of vitality: assembly theory, Tononi's integrated information, Deacon's teleodynamics, Friston's FEP. The comparative diagnosis: each works at the level of \(I_v\) or of a single capacity; none reaches \(\eta_L\) as a criterion of a closed self-payment loop.

\subsection{Assembly Theory}\label{assembly-theory}

Marshall, Murray, Cronin \cite{Marshall2017} formalized the assembly index \(a\) --- the minimal number of recursive assembly steps of an object with the reuse of parts, measured experimentally through the mass spectrometry of complex molecules; the extended program was formulated by Walker and coauthors \cite{Walker2023}. An object with \(a\) above an empirical threshold (of the order of 15 for biogenic molecules) is interpreted as a biosignature: the depth of accumulated selection required to assemble such an object is not attained by abiotic processes on observable timescales.

The assembly index works on the denominator of the ``model × self-payment'' dyad --- a measure of accumulated selection --- but does not define the numerator. Assembly theory assigns a high \(a\) equally to a fossil and to a living cell, failing to distinguish the selection accumulated in matter from the active maintenance of \(I_{\text{pred}}\); the coefficient of error return in the object's current behavior does not enter the scheme. On the \(\eta_L\) scale the assembly index is an operational complexity index, not a vitality scale: for systems with a complex molecular architecture a high \(a\) is a necessary condition for a positive \(\eta_L\), but not a sufficient one. The assembly index is an operational proxy for \(C_v\) or \(S_v\) within \(I_v\), not a self-standing scale.

\subsection{\texorpdfstring{Integrated Information \(\Phi\)}{Integrated Information Phi}}\label{integrated-information-phi}

Tononi and his school \cite{Tononi2004,Tononi2016} defined integrated information \(\Phi\) as a measure of the non-decomposability of a whole into a sum of parts. In the formalization of IIT 3.0 \cite{Oizumi2014}, \(\Phi > 0\) is interpreted as a sufficient condition for minimal phenomenal consciousness. In the context of the demarcation of the vital, \(\Phi\) is a powerful candidate for the formalization of the informational capacity \(C_v\): systems with high integrated information satisfy the requirement of a model in which the parts are not decomposable into independent subsystems.

Tononi's program imposes no constraint on the thermodynamic payment for maintaining this integrated information and does not require that the payment belong to the system itself. On the weak formulations of IIT a household thermostat has a nonzero \(\Phi\) --- hence the conclusions discussed in \cite{Tononi2016} about the minimal ``experience'' of simple systems. \(\Phi\) is a necessary condition for a positive \(C_v\), not a sufficient condition for a positive \(\eta_L\): thermodynamic ownership of the payment by the system remains outside the IIT apparatus. The connection is productive as a decomposition: \(\Phi\) operationalizes \(C_v\), and \(\eta_L\) adds the requirement \(T_v > 0\) with belonging to the same system.

\subsection{Teleodynamics}\label{teleodynamics}

In \emph{Incomplete Nature} \cite{Deacon2011} Deacon built a hierarchy homeodynamics → morphodynamics → teleodynamics. The third stage is defined as the coupling of two morphodynamic systems that mutually constrain each other's self-destruction; the notion of \emph{absential causality} is introduced --- causation by that which is not yet (by a virtual goal, not given actually). Teleodynamics is the program closest to the ``model × self-payment'' dyad: the model and the self-payment are present in the scheme logically, both axes are articulated.

The defect of the program is the absence of an operational quantitative measure. Deacon's hierarchy is stated in qualitative philosophical terms, and reconstructing \(\eta_L\) or a coefficient of error return through absential causality is technically impossible without a substantial formalization that the program does not offer. Teleodynamics is a structural portrait in which both axes of the dyad are present, but a quantitative measure, testable on a bacterium or a metropolis, is absent. Teleodynamics is a conceptual predecessor whose operationalization through the \(\eta_L\) scale is precisely one of the contributions of the present work.

\subsection{The Free Energy Principle}\label{the-free-energy-principle}

Friston's free energy principle \cite{Friston2010} is the formalization closest to the present framework. In the FEP literature two quantities are distinguished. The \emph{variational free energy} (VFE)

\begin{equation}
F[q, o] = D_{KL}\bigl[q(s) \,\|\, p(s|o)\bigr] - \ln p(o), \tag{3}
\end{equation}

is minimized in perception: the distribution \(q(s)\) over the hidden states of the environment approximates the true posterior \(p(s|o)\); \(-F\) is a lower bound on the log-evidence \(\ln p(o)\). The \emph{expected free energy} (EFE)

\begin{equation}
G[\pi] = \mathbb{E}_{q(o', s'|\pi)}\bigl[\ln q(s'|\pi) - \ln p(o', s')\bigr], \tag{4}
\end{equation}

is minimized in action selection / active inference. Here \(p(o', s')\) is a preference distribution (a generative model with preferences over future outcomes and states), and \(G\) canonically decomposes into epistemic value (the informational gain about hidden states) and pragmatic value (closeness to preferences). The difference is fundamental: \(F\) describes the updating of beliefs, whereas \(G\) governs the selection of actions \cite{Friston2010,Friston2019}. For a Gaussian generative model in a stationary regime, \(-F \lesssim I_{\text{pred}}\) --- a heuristic relation linking the scale and FEP in the perceptual part; a rigorous derivation with non-stationary corrections is an open technical problem.

The principled objection to FEP turns on disentangling two interpretations of the Markov blanket. Bruineberg, Kiverstein, and Rietveld \cite{Bruineberg2018} give an ecological-enactive critique of pure FEP, and Bruineberg, Dołęga, Dewhurst, and Baltieri \cite{Bruineberg2022} give a mature formalization of the distinction between the \emph{Pearl blanket} and the \emph{Friston blanket}. The Pearl blanket is an epistemic instrument of Bayesian inference about a system from outside; the Friston blanket is an ontological property of the system's stationary density. The authors showed that pure FEP has no resources for the second interpretation without additional ontological premises, illustrating this with the household thermostat, which formally satisfies FEP but is neither a living nor a mental system. The thermostat illustration in Bruineberg et al.~is auxiliary; the central argument is precisely this disentangling of the Pearl and Friston blankets.

Our move --- the introduction of the self-payment requirement --- adds an ontological premise not derivable from pure FEP: \(F\) pertains to the system's own dissipation. Functionally this is analogous to the transition from the Pearl- to the Friston-blanket through a thermodynamic channel, not a statistical one. For a thermostat the predictive information about its own environment, held by the internal circuit, is negligible: the setpoint is imposed externally, and the circuit has no own model of the environment paid for by its own dissipation. The numerator of the \(\eta_L\) of a thermostat tends to zero; the Pearl-blanket remains applicable, the Friston-blanket does not.

\emph{Counterfactual screening of the thermostat.} Apply the procedure of § 2.2 to a bimetallic thermostat with hysteresis \(\Delta T_{\text{hyst}} = 1\) K, controlling a room with thermal inertia \(\tau_{\text{room}} \sim 30\) min. Five candidates for inclusion in the boundary are sifted: the bimetallic strip (sensor + actuator), the relay contact group (energy gating), the heating element, the power grid as a current source, and the external setter. Counterfactual subtraction leaves the first two inside the boundary, since removing either breaks the regulation circuit. The other three are external: subtracting the heater removes the thermal source, not the decision-making loop; subtracting the power grid is analogous; and subtracting the setter does not disrupt regulation around a fixed setpoint. Inside the boundary, \(E_{\text{actual}}\) is the elastic relaxation of the metal plus the dissipation of the relay, \(\sim 10^{-6}\) W (a reproducible computation of the two interpretations --- trivial versus managing --- is in the Supplementary Materials). The internal channel that holds \(I_{\text{pred}}\) about the room temperature is the deformation of the bimetal, with \(I_{\text{pred}} \le \log_2(\Delta T_{\text{room}}/\Delta T_{\text{hyst}}) \sim 5\) bits. This \(I_{\text{pred}}\) pertains to the environment ``room,'' not to the environment ``thermostat + room + power source.'' The loop of decisions paid for by the system's own breakdown is therefore incomplete for the thermostat: the room is heated by the external power grid, not by the bimetal's own free energy. \(\eta_L^{\text{thermostat}} > 0\) only for the trivial task of ``maintaining the bimetal's own shape,'' not for the managing task of ``regulating the room temperature.'' The divergence between the two interpretations is exactly the transition from the Pearl to the Friston blanket: one and the same physical system admits several self-consistent boundaries, and the choice between them is a methodological discipline. The full per-component analysis of the five candidates is in Supplementary § S4.4.

The relation between the classes of FEP-applicable and \(\eta_L\)-vital systems is asymmetric. The claim that ``FEP is a special case of \(\eta_L\)'' has the direction wrong. The correct relation is the reverse: within the class of FEP-applicable systems, \(\eta_L\) isolates a subclass in which the system's own dissipation pays for the minimized \(F\). Outside this subclass --- a thermostat, a pendulum with external power, any self-regulating circuit with externalized payment --- \(\eta_L\) goes to zero through the numerator or the denominator. FEP retains its descriptive power over the entire class but loses its discriminating power: it does not separate systems with self-payment from systems without self-payment. \(\eta_L\) is not an alternative to FEP but its thermodynamic discipline: an operationalization that restores discriminating power by an explicit indication of ownership.

\emph{Active inference and the space of correspondences.} Active inference \cite{Friston2017,Parr2022} is a concrete computational scheme derivable from FEP under policy selection via EFE. In the terms of active inference the numerator of \(\eta_L\) is connected with predictive information in the sense of a generative model, and the denominator with the pragmatic value of policies through a thermodynamic channel. An extended comparison of the vocabularies (action, policy, pragmatic value, epistemic value vs.~\(I_{\text{pred}}\), \(E_{\text{actual}}\), \(T_v\), \(S_v\)) is a task for a separate work; what is essential here is that the correspondence is technically possible and does not require revising any of the source frameworks.

\emph{Andrews 2021 and Williams 2020.} Andrews \cite{Andrews2021} attacks FEP for the non-derivability of empirical content from the formalism; the critique is parried in \(\eta_L\) by tying the numerator and the denominator to measurable quantities. Williams \cite{Williams2020} attacks predictive coding as a neural mechanism, not predictive information as a formal measure; the critique is neutral to \(\eta_L\). An extended analysis and the distinction between predictive coding \cite{RaoBallard1999,Clark2013} vs.~predictive information \cite{Bialek2001} are in Supplementary § S4.4a.

Summary position: the formalism of free-energy minimization becomes a substantive criterion of vitality only under the condition of self-payment. Without this condition FEP describes any self-regulating circuit and loses its discriminating power. With this condition FEP is refined to its \(\eta_L\)-discipline: for systems where FEP is applicable in the strict sense, \(\eta_L\) and FEP give consistent characterizations; for systems where FEP is applicable only in the weak formulation, \(\eta_L\) goes to zero through the numerator.

\subsection{Demarcation from Adjacent Programs: A Summary Diagnosis}\label{demarcation-from-adjacent-programs-a-summary-diagnosis}

The summary comparative diagnosis of the four programs is collected in the table below and lays bare a single structural pattern of the defect: each program takes exactly one side of vitality seriously and leaves the rest outside its apparatus.

\begin{longtable}{@{}p{\dimexpr 0.1500\linewidth-2\tabcolsep\relax}p{\dimexpr 0.3667\linewidth-2\tabcolsep\relax}p{\dimexpr 0.4833\linewidth-2\tabcolsep\relax}@{}}
\toprule\noalign{}
\begin{minipage}[b]{\linewidth}\raggedright
Program
\end{minipage} & \begin{minipage}[b]{\linewidth}\raggedright
Axis taken seriously
\end{minipage} & \begin{minipage}[b]{\linewidth}\raggedright
Left outside the apparatus
\end{minipage} \\
\midrule\noalign{}
\endhead
\bottomrule\noalign{}
\endlastfoot
Walker--Cronin (assembly theory) & the denominator of the dyad --- accumulated selection in matter (assembly index \(a\)) & the numerator is implicit; \(a\) is equal for a fossil and a living cell --- vital discrimination requires an external intuition about the error-return loop \\
Tononi (IIT) & a measure of informational integration \(\Phi\) & the thermodynamic payment; the \(\Phi\)-positive thermostat is a known critical point of the theory \\
Deacon (teleodynamics) & a structural dyad (absential causality) & a quantitative measure: absential causality is not reconstructed through measurable quantities \\
Friston (FEP) & a formalism of free-energy minimization & the requirement that model and dissipation belong to one system; disentangling the Pearl and Friston blankets in Bruineberg et al.~fixes precisely this gap (see § 4.4) \\
\end{longtable}

The four defects are not accidental: they reflect four ways of taking one side of vitality seriously while leaving the rest outside the apparatus. As a ratio that requires both components to belong to one system, \(\eta_L\) is an operationalization that lets none of the four axes dissolve into the background.

\subsection{Connection to the Tradition: Six Lines of Convergence}\label{connection-to-the-tradition-six-lines-of-convergence}

\(\eta_L\) converges at a single point with six classical lines. Schrödinger's ``life feeds on negentropy'' \cite{Schrodinger1944} anticipates the Landauer denominator \(N_{\text{max}}\) as a formalization of the price of order. Wiener \cite{Wiener1948} supplies the structural apparatus of the loop ``compression → self-payment → correction'' but without requiring that it belong to one system. Rosen's anticipatory systems \cite{Rosen1985} are a conceptual predecessor of self-modeling without a thermodynamic operationalization. Taleb's skin in the game \cite{Taleb2018} is the concept of self-payment, here formalized thermodynamically in § 2.1. Kolmogorov's algorithmic information and induction through compression are a predecessor of \(C_v\). Kauffman's autocatalytic closures \cite{Kauffman1993,Kauffman2000} --- ``order for free'' plus a thermodynamic working cycle --- supply half of the denominator without the numerator.

The section's summary comparative conclusion is as follows. Walker--Cronin gives the quantity of accumulated selection without a coefficient of error return. Tononi gives a measure of informational integration without requiring the system's own thermodynamic payment. Deacon gives the structural dyad of model and self-payment without a quantitative measure. Friston gives a formalism of free-energy minimization without requiring that this energy belong to the modeling circuit. Schrödinger, Wiener, Rosen, Taleb, Kolmogorov, and Kauffman are conceptual predecessors, each working on one axis of the dyad. As a ratio, the \(\eta_L\) scale assembles both axes of the dyad into one quantity and explicitly introduces the requirement of the system's own dissipation. This does not deny the listed programs but indicates their comparative place: each is operationalized through one of the capacities within \(I_v\), and unifying them into one vitality scale is a self-standing step that none of the source frameworks took.

\section{Falsification Criteria}\label{falsification-criteria}

A procedural rule common to all lines below: the system boundary and the proxies for \(I_{\text{pred}}\) and \(E_{\text{actual}}\) are fixed before the \(\eta_L\) verification and are not revised to obtain the desired sign. Any revision of the boundary is a separate paradigm-case counterfactual-screening procedure (§ 2.2), not an ad hoc rescue tailored to the result.

The program admits falsification along eight independent lines, grouped into three blocks by the character of the evidence.

\emph{Direct refutations of the scale.} The first three criteria attack the structure of \(\eta_L\) through the zeroing of one of the three capacities.

\begin{enumerate}
\def\labelenumi{\arabic{enumi}.}
\tightlist
\item
  The discovery of a system that maintains \(I_{\text{pred}} > 0\) stationarily at an \(E_{\text{actual}}\) below the Landauer bound \(I_{\text{pred}} \cdot k_B T \ln 2\) in an irreversible regime of realization. Such an observation would violate Landauer's principle on the class of systems for which \(\eta_L \le 1\) --- the Proposition (§ 2.1) --- and would zero \(T_v\) as a necessary condition. The Bennett regime of reversible computation \cite{Bennett1973,Bennett1982} does not count as a refuting scenario: it lies outside the domain of definition and requires a separate operationalization through the number of logical bits.
\item
  The discovery of a system with a cumulative complexification of its own model of the environment without stochastic variation, without errors, without a correction loop --- a structure that evolves deterministically while adapting. This would zero \(C_v\).
\item
  The discovery of a system with stable self-paying compression under a fundamentally non-hierarchical, non-modular topology of connections --- in an Erdős--Rényi random graph without modular structure or in a strictly regular lattice. This would zero \(S_v\).
\end{enumerate}

Refuting any capacity simultaneously collapses both scales: \(\eta_L\) --- through the numerator or the denominator, \(I_v\) --- through the geometric mean. This is one falsification through a capacity, not a double one. The program does not provide for an ad hoc rescue of the form ``one scale remains.''

\emph{A condition for confirming the uniqueness of the operationalization (not a falsification criterion).} The discovery of an alternative proxy that gives the same demarcation discrimination but is not reducible to \(\eta_L\) either through identity or through a monotone transformation does not refute the program. It confirms the substantiveness of the demarcation and raises the question of competition among equivalent scales. Symmetrically, the absence of such a proxy over long times is indirect evidence of the uniqueness of \(\eta_L\).

\emph{Tests of the proxies and the operationalization.} The next three criteria attack not the structure of the scale but its tie to measurable quantities.

\begin{enumerate}
\def\labelenumi{\arabic{enumi}.}
\tightlist
\item
  A systematic violation of the monotone dependence of the rate of evolutionary learning on \(\eta_L\). If, in a controlled experimental-evolution experiment (bacterial cultures in a chemostat under the variation of temperature, population density, and carbon source), systems with \(\eta_L\) differing by an order of magnitude give rates of adaptation that are statistically indistinguishable over sufficiently long windows, the central ratio loses its numerical meaning. The constructive check coincides with the protocol of the ``Progressive predictions'' block below.
\item
  The discovery of a system that indicates vitality through three independent channels --- the rate of adaptation, the rate of self-repair, and independent biochemical markers of self-reproduction (in the sense of Darwinian evolution; cf.~the NASA definition of life as ``a chemical system capable of Darwinian evolution'' \cite{Joyce1994}, invoked here as a conventional external referent, not as a competing theory of vitality, the deliberate switch of the substrate-bound register being intentional) --- while \(\eta_L = 0\) under all operationalizations of \(I_{\text{pred}}\) and \(E_{\text{actual}}\) through standard domain-specific proxies. This criterion is independent of \(\eta_L\): the rate of adaptation, the rate of self-repair, and the biochemical markers of self-reproduction are operationalized through measurable quantities not reducible to the numerator of \(\eta_L\). A systematic divergence would point either to a defect of the numerator as an operational proxy for predictive information (the methodological level of falsification, see § 2.7) or to the untenability of the very concept of vitality through the efficiency of self-modeling (the ontological level).
\item
  The discovery of a biological system, unquestionably alive by conventional criteria, for which \(\eta_L\) is strictly equal to zero under all reasonable operationalizations.
\end{enumerate}

\emph{Progressive predictions.} The final block fixes the excess empirical consequences of the program, not derivable from any of the competing frameworks.

\begin{enumerate}
\def\labelenumi{\arabic{enumi}.}
\tightlist
\item
  \emph{A quantitative prediction about the rate of evolutionary learning.} The rate of decrease of the divergence between the internal model and the environment, normalized to the complexity of the environment, must scale with \(\eta_L\) as a monotonically increasing function. The concrete protocol uses a culture of \emph{E. coli} in a chemostat with a controlled carbon flow (glucose, lactose, acetate). \(\eta_L\) is varied through temperature (25--42 °C, five levels) and population density (three levels), and adaptation is the decrease of the divergence between the encoded methylation level and the observed gradient over 50 generations. The positive prediction is \(\rho \in [0.4,\, 0.7]\) on 45 lines and \(r_{\text{adapt}} = A \cdot \eta_L^{\alpha}\) with \(\alpha \in [0.3,\, 1.0]\). The upper bound \(\alpha = 1\) is a direct proportionality of the rate of adaptation to the flux of \(I_{\text{pred}}\); the lower bound \(\alpha = 1/3\) follows from the structure of (2), since \(I_v\) is the cube root of the product of three capacities and the minimal dependence \(r_{\text{adapt}} \propto \sqrt[3]{\eta_L}\) corresponds to the contribution of one capacity out of three. The scale is calibrated on the worked example of § 3.5 (\(\eta_L^{\text{comp}} \in [10^{-6}, 10^{-4}]\)). Refutation: \(\hat{\alpha}\) statistically distinguishable from \([0.3,\, 1.0]\) at \(p < 0.01\); \(\rho < 0.2\) at the stated power; or \(\rho < 0\).
\item
  \emph{A summary Lakatosian passport of the program.} Relative to assembly theory, IIT, teleodynamics, and FEP the program is a progressive shift on three predictions, following from none of the source frameworks: the monotone correlation of the rate of adaptation with \(\eta_L\) (see above); a transient pattern of SETI signals as a signature of the payment-externalization filter (§ 6); the abrupt transition of an autonomous embodied agent into vital status upon crossing the loop-closing threshold. The program sustains the progressive regime as long as these predictions admit independent empirical testing.
\end{enumerate}

\section{Discussion and Open Problems}\label{discussion-and-open-problems}

\emph{Astrobiology.} The principal consequence is a reformulation of the search for life. Most objects pass the \(I_v\) screening (the structural potential is universally high) and fail the \(\eta_L\) verification (actual closure of the self-payment loop is rare). Operationally, when designing biosignature search missions (Europa Clipper, JUICE, Dragonfly, exoplanet spectroscopy), the priority shifts from chemical disequilibrium to the search for a triple simultaneity: non-equilibrium chemistry (\(T_v\)), a biogenic-like informational signature (\(C_v\), through isotopic fractionation and chirality), and a hierarchically modular topology of the environment (\(S_v\)). Each is a necessary condition; the simultaneity of all three is an operational candidate for a biosignature.

\emph{The ethics of artificial intelligence.} The question shifts from ``is the model complex enough to be a moral subject?'' to ``will the model itself begin to physically bear the cost of its own errors?'' \(\eta_L\) turns from zero to positive precisely when an own dissipative payment appears for a model held in the working circuit's own degrees of freedom. The current LLM architecture excludes this transition by construction; this is a structural condition, not a calendar forecast.

``Bearing its own errors'' is meant in the technical sense of § 2.1: the model's errors return as a physical loss of the system's own free energy, not as reputational sanctions through an external corporate substrate. A fall in an LLM's usage rate, with retraining by the corporation, is the corporation's payment and enters the corporation's \(\eta_L\); the \(\eta_L^{\text{int}}\) of the model as an agent stays zero. Regulatory fines, legal sanctions against the developer, and market reputational losses return harm not to the system that holds the model but to its external substrate.

\emph{Complex systems and social dynamics.} In application to cities and civilizations the scale opens two thresholds: structural screening (\(I_v\)) and verification (\(\eta_L\)). A civilization with a high \(I_v\) and a sagging \(\eta_L\) --- structurally rich, failing the maintenance of the loop --- is diagnostically distinct from a civilization with a sagging screening (the destruction of infrastructure). The differentiation is useful for research on the resilience of social systems; it requires empirical verification.

\emph{The Fermi paradox and the externalization filter.} In the terms of the framework a civilization passes the ``filter'' if and only if its \(\eta_L\) is positive with the numerator belonging to its own dissipation and its \(T_v\) held inside its own circuit. Most hypothetical candidates may sag through a rupture of self-payment: they pass the \(I_v\) screening (radio signals, megastructures) but fail the \(\eta_L\) verification through the externalization of payment. This claim has the status of a speculative extrapolation and is subject to independent empirical testing. The testable statistical prediction about observable SETI signals is twofold: transient, non-repeating signals with a power-law duration distribution and a cutoff indicate the filter, whereas a stationary spectrum with an IR excess holding over decade-long windows indicates passage of both thresholds. The refutation condition: the discovery of three or more independent stationary SETI sources with a confirmed IR signature of megastructures over intervals \(\gtrsim 50\) years is incompatible with the prediction of the externalization filter in its current formulation.

\emph{The status of the consequence.} The reformulation has the status of a speculative extrapolation: neither the NASA Astrobiology Strategy, nor the ESA Cosmic Vision, nor the planetary Decadal Survey operates with the category of a ``self-payment loop.'' If the standard programs attain a reliable demarcation of extraterrestrial life without recourse to \(\eta_L\), the framework will bring no astrobiologically new discriminating power and reduces to the existing agnostic biosignatures.

\emph{Open problems.} A non-stationary version of the stationary-erasure lemma (§ 2.1) for systems with growing memory --- the biosphere on eonic windows, learning artificial systems --- requires a rigorous proof that accounts for the cost of memory expansion. The catalogue of operational proxies for \(I_{\text{pred}}\) must be expanded with numerical examples and limits of applicability. The case of Gaia --- the unity of the biosphere's model --- is empirically open and depends on data about the coherence of biogeochemical cycles over eonic windows. The operational protocol for detecting a vital AI requires precise thresholds for population experiments. The philosophical status of \(\eta_L\) --- a position of moderate operationalism with a demarcation of the ontological referent and the operational measure --- is a task for a separate work.

\emph{Boundary conditions of applicability.} The \(\eta_L\) scale does not make three types of claims.

\begin{enumerate}
\def\labelenumi{\arabic{enumi}.}
\tightlist
\item
  It does not assert the identity of vitality and consciousness. A high \(\eta_L\) is a necessary condition for complex cognitive activity (without the maintenance of one's own predictive information neither prediction, nor planning, nor reflection is possible), but not a sufficient one. The question of phenomenal experience, discussed within the framework of integrated information \cite{Tononi2016} and the hard problem of consciousness, lies beyond the \(\eta_L\) scale.
\item
  It does not presuppose biocentrism by a reverse gesture. A high \(\eta_L\) is possible for systems of a non-biological substrate --- major cities, planetary-scale biospheres, hypothetical artificial embodied agents with a closed self-payment loop. This is a substantive prediction, not a calibrational fitting.
\item
  It does not reduce to a moral criterion. A civilization that fails the \(\eta_L\) verification with the \(I_v\) screening preserved is diagnosed as structurally rich but non-living --- a diagnosis of a structural state, not an ethical verdict. The ethical consequences of the program are a separate topic, resting on the structural diagnosis but not derivable from it by direct implication.
\end{enumerate}

\section{Conclusions}\label{conclusions}

The work has defined vitality operationally, as the Landauer efficiency of self-modeling: the dimensionless ratio \(\eta_L = I_{\text{pred}} / N_{\text{max}}\) of the predictive information the system holds about its own environment to the Landauer budget of its own dissipation in the irreversible regime. Positivity requires the simultaneous presence of three capacities --- \(T_v\), \(C_v\), \(S_v\) --- assembled into a composite index of structural prerequisites \(I_v\). The index is necessary but not sufficient for vitality and occupies the same comparative place as the assembly index, \(\Phi\), and Deacon's morphodynamics. \(\eta_L\) itself is the only logically substantive criterion of vitality, in the sense of a binary sign (\(\eta_L > 0\) with a closed self-payment loop), and it distinguishes structurally rich but non-living systems (the Sun, a hurricane, an LLM under the boundary ``weights at the moment of inference'') from systems with \(\eta_L > 0\) (a bacterium, a city, the biosphere). What is substantive across classes is precisely the sign, not the absolute magnitude (§ 3.5).

The application of the two-stage procedure to six paradigm cases yields sharp demarcation answers and lays bare a structural asymmetry: complexity is ubiquitous, closed self-payment loops are rare. Comparison with assembly theory, IIT, teleodynamics, and FEP positions \(\eta_L\) as an operational closure of the gap of each of the programs. Concrete falsification conditions are formulated, tied to testable physical, informational, and evolutionary predictions.

The program is open to refutation along eight independent lines (with a separately fixed condition for confirming the uniqueness of the operationalization) and is closed to ad hoc rescues by its two-stage architecture: refuting one capacity collapses both scales. Herein lies its scientific status --- not a claim to a final definition of life but a testable research program with explicitly articulated limits of applicability.

The principal consequence of the shift from ``vitality as substance'' to ``vitality as the efficiency of self-modeling'' is an operational reformulation of traditionally metaphysical questions. The search for extraterrestrial life becomes a search for closed self-payment loops, not for complex chemistry. The ethical status of artificial systems becomes a question of the appearance of an own dissipative payment for the model, not of functional similarity to a human. The boundary between the living and the non-living becomes operational: it runs where \(\eta_L\) ceases to be zero. Each of these reformulations is not an automatic solution of the corresponding question; it is an indication of a concrete measurable quantity with respect to which the question can be formulated in a disciplined way.


\vspace{6pt}

\supplementary{The following supporting information is deposited with this article: (i)~Supplementary Notes (extended operationalization protocols, full per-component counterfactual analyses, sensitivity discussion); (ii)~reproducible Python code for the worked-example calculations of Sections~3.4--3.7 and~4.4 (\textit{E.~coli} chemotaxis, the bimetallic thermostat, a major city, the biosphere under the single-compressor reading, and LLM-as-corporation). Each script uses only the Python standard library and produces a captured \texttt{expected\_output.txt} for verification. A versioned archive is deposited on Zenodo: \url{https://doi.org/10.5281/zenodo.20262947}.}

\authorcontributions{Conceptualization, methodology, formal analysis, writing---original draft preparation, writing---review and editing: A.A. The author has read and agreed to the published version of the manuscript.}

\funding{This research received no external funding.}

\institutionalreview{Not applicable.}

\informedconsent{Not applicable.}

\dataavailability{No primary empirical data were generated for this work. The reproducible computations underlying the worked examples of Sections~3.4--3.7 and~4.4 are provided as part of the Supplementary Materials (Python source code, captured outputs, parameter documentation). The open-access materials---the manuscript, supplementary material, bibliography, graphical abstract, and the reproducible worked-example code---are publicly available in the GitHub repository \url{https://github.com/andriishin/landauer-self-modeling-oa} and in a versioned Zenodo archive: \url{https://doi.org/10.5281/zenodo.20262947}.}

\vspace{6pt}\noindent{\fontsize{9}{11.2}\selectfont\textbf{Use of AI Tools:} The author used Anthropic Claude (Opus 4.6, 4.7) as an editorial and infrastructure assistant during the preparation of this manuscript and its supplementary materials. The tool was used for: (i)~translation and iterative revision of the English text (originally drafted in Russian) against structured checklists derived from the journal's author guidelines (length corridor, structural completeness, epistemological-marker discipline, falsifiability framing); (ii)~generation and verification of the reproducible Python code in the Supplementary Materials; (iii)~cross-checking of bibliographic metadata against locally cached copies of the cited sources; and (iv)~operation of the \LaTeX{} build pipeline. All scientific claims, the $\eta_L$ formulation, the self-payment requirement, paradigm-case diagnoses, and falsification criteria were formulated by the author. AI tools did not generate scientific content autonomously. The author reviewed every AI-assisted output and takes full responsibility for the content of this manuscript and its supplementary materials.\par}

\conflictsofinterest{The author declares no conflict of interest.}

\reftitle{References}

\externalbibliography{yes}
\bibliography{refs}

\end{document}
